\documentclass[11pt]{article}

\usepackage[margin=1in]{geometry}
\usepackage[T1]{fontenc}
\usepackage{csquotes}
\usepackage{graphicx}
\usepackage{booktabs}

\usepackage[backend=biber,style=numeric,sorting=none,maxnames=3,minnames=1]{biblatex}
\addbibresource{bib/refs.bib}

\usepackage{hyperref}

\usepackage{bucket}

\hypersetup{
  colorlinks = true,
  linkcolor  = bucketblue,
  citecolor  = bucketblue,
  urlcolor   = bucketblue,
}

\newcommand{\emp}[1]{{\small[empirical: #1]}}
\newcommand{\bm}[1]{{\small[bm:#1]}}
\newcommand{\sumrun}{data/summary.json, 2026-11-26, python3 scripts/stats.py}
\newcommand{\recrun}{data/recall.json, 2026-11-26, python3 scripts/recall.py}
\newcommand{\testrun}{data/tests.json, 2026-11-26, python3 scripts/tests.py}

\glossentry{forest}{inheritance forest}%
  {https://github.com/gianyrox/bucket-foundation/blob/dev/lean/BucketMath/Etymology/Forest.lean}%
  {The graph whose nodes are lexemes and reconstructed roots and whose
   edges point from each node to the one form it was inherited from; see
   \Cref{def:forest}.}
\glossentry{rootdist}{root distance}%
  {https://github.com/gianyrox/bucket-foundation/blob/dev/papers/root-graph/scripts/tests.py}%
  {The number of inheritance steps from two lexemes up to their lowest
   common ancestor, summed; see \Cref{def:rootdist}.}
\glossentry{bshare}{borrowed share}%
  {https://github.com/gianyrox/bucket-foundation/blob/dev/papers/root-graph/PREREGISTRATION.md}%
  {Borrowed words divided by borrowed plus inherited words, among words
   whose first etymology relation is one of the two.}

\title{A Root Graph for Every Language in Wiktionary}
\author{Bucket Foundation}
\date{2026-11-26}

\begin{document}
\maketitle

\begin{abstract}
We build an etymology graph over every language in English Wiktionary:
6,970,532 nodes for 5,967 language codes, 158,107 of them reconstructed
forms, and 6.5 million edges of types inherited, borrowed, derived,
cognate, calque and affix \emp{\sumrun}. Descent edges, inheritance plus
derivation from a language on the target's ancestral line, form a
\term{forest}{forest} with 370 family roots and depth up to 15. On 90
hand-listed cognate pairs the forest links 87, against 75 in a first
build without ancestral-line derivation \emp{\recrun}. We prove in Lean
that the forest is acyclic, that lowest-common-ancestor depth is an
ultrametric similarity, that the path distance we test with is a metric
and fails the ultrametric inequality, and that borrowing breaks the
ultrametric only through a word with two ancestry paths
\bm{BucketMath.Etymology.single\_address\_ultra}. Three preregistered
tests follow, run after a dated amendment to the graph.
\term{rootdist}{Root distance} predicts the phonetic similarity of
translation pairs, Spearman $\rho=-0.27$, 95\% CI $[-0.36,-0.18]$,
$n=464$, Holm $p=4.7\times10^{-9}$. We fail to reject H0 for semantic
similarity, $n=80$. Frequent vocabulary has a lower
\term{bshare}{borrowed share} than the whole lexicon in 193 of 206
languages, 158 significant after Benjamini-Hochberg. We fail to reject H0
for borrowed share against cognate-coverage residuals, $\rho=0.26$, CI
$[-0.18,0.62]$, $n=23$: inconclusive \emp{\testrun}.
\end{abstract}

\section{Introduction}
\label{sec:intro}

A learner who knows that German \emph{Wasser} and English \emph{water}
come from one Proto-Germanic form has a reason to expect the two to sound
alike. A learner who meets Japanese \emph{tēburu} has a different reason:
Japanese borrowed it from English. Course sequencing in Bucket's learning
engine uses a tree similarity, the depth of the shared prefix of two
taxonomy addresses, and the proof that every acquisition order costs the
same rests on that similarity being an ultrametric
\bm{BucketMath.Learning.Sequencing.lcp\_isUltra}. Etymology gives a
natural tree for words. Borrowing makes it a graph. We build the graph at
full Wiktionary scale, prove what survives of the tree structure, and test
three empirical claims against it.

\section{Related Work}
\label{sec:related}

Wiktextract turns Wiktionary pages into structured records, including the
etymology templates we parse \autocite{ylonen2022wiktextract,kaikki2026raw}.
The World Loanword Database measured borrowing rates over 41 languages and
found core vocabulary resists borrowing \autocite{haspelmath2009loanwords};
our T2 repeats that question over 198 languages with a cruder
classification. Phonetic alignment for cognate detection is surveyed by
List \autocite{list2014sequence}. Ultrametrics on rooted trees are
standard in phylogenetics \autocite{semple2003phylogenetics}.

\section{Graph}
\label{sec:graph}

\begin{definition}[Root graph]
\label{def:graph}
Nodes are pairs (language code, entry form). Etymology-only varieties,
such as Late Latin or Old Italian, share the code of the language they
belong to in Wiktionary's language data. Forms are cleaned: HTML entities
unescaped, wiki brackets and inline modifiers removed, comma-joined forms
cut to the first, \emph{oe} and \emph{ae} ligatures spelled out, NFC, and
Wiktionary's entry-name rules applied (macrons and breves everywhere,
accents in Slavic languages and Proto-Slavic, overdots in Old and Middle
English). A form whose every sense is a form-of one lemma, and which has
no lemma entry of its own, maps to that lemma, so Latin \emph{lactem}
joins \emph{lac}. A node is reconstructed when its form starts with an
asterisk or its language is a proto-language. Edges come from etymology
templates and from the ancestry trees that the \texttt{etymon} template
prints: inherited, borrowed, derived, calque, cognate and affix, each with
the source page and template name.
\end{definition}

\begin{definition}[Ancestral line]
\label{def:line}
The ancestral line of a language is the transitive closure of the
\texttt{ancestors} fields in Wiktionary's language modules, the parents of
etymology-only varieties, the proto-language of the nearest family when a
language names no ancestor, and inheritance language pairs seen at least
20 times in etymon trees.
\end{definition}

\begin{definition}[Inheritance forest]
\label{def:forest}
A descent edge is an inheritance edge to a language on the ancestral line,
or a derivation edge from one. Each node keeps its first descent claim in
file order, an inheritance claim before a derivation claim on the same
page. Claims that would close a cycle are dropped. The result is a
forest, \Cref{thm:acyclic}.
\end{definition}

Borrowing, learned borrowing and calque stay outside the forest. A
derivation from a language off the ancestral line is its own class,
\texttt{der\_ext}. The build drops 3,835 inheritance claims that point off
the line, keeps 126,481 derivation edges as descent and drops 187
cycle-closing claims \emp{\sumrun}. It streams the 3.1 GB compressed
dump twice, once for the form-of map and once for edges, in one process
under a 20 GB memory cap, in eleven minutes.

\begin{table}[htbp]
  \centering
  \caption{Graph size \emp{\sumrun}.}
  \label{tab:graph}
  \begin{tabular}{@{}lr@{}}
    \toprule
    Item & Count \\
    \midrule
    Nodes & 6,970,532 \\
    Reconstructed nodes & 158,107 \\
    Language codes with entries & 5,162 \\
    Affix edges & 3,252,206 \\
    Inherited edges & 1,079,346 \\
    Derived edges & 842,406 \\
    Cognate edges & 733,497 \\
    Borrowed edges & 574,094 \\
    Calque edges & 29,628 \\
    Family roots with descendants & 370 \\
    \bottomrule
  \end{tabular}
\end{table}

The largest families by rooted lexemes are Proto-Indo-European (113,380
lexemes over 480 languages), Latin (46,474), Proto-Slavic (26,970),
Proto-Germanic (23,511) and Proto-Turkic (17,546). Latin and Middle
English appear as roots where Wiktionary stops the chain there. Per-language
and per-family tables are in \texttt{\_intake/root-graph/stats\_lang.csv}
and \texttt{stats\_family.csv}.

\paragraph{Linkage recall.} \texttt{data/cognates.tsv} lists 90 cognate
pairs, six pairs for each of five concepts within Spanish, Italian,
Portuguese and French; English, German, Dutch and Swedish; and Russian,
Polish, Czech and Ukrainian. A pair counts as linked when the two forest
chains meet. The second build links 87 of 90: 27 Romance, 30 Germanic, 30
Slavic. The first linked 75, with 21 Slavic. The three misses all involve
French \emph{lait}, which Wiktionary derives from a Vulgar Latin form
distinct from the Latin lemma. Over all translation-table pairs within
each group, non-cognate synonyms included, the linked share rises from
1.1\% to 2.6\% in Romance, 0.9\% to 1.2\% in Germanic and 1.1\% to 1.7\%
in Slavic \emp{\recrun}.

\section{Proofs}
\label{sec:proofs}

\begin{theorem}[Acyclicity]
\label{thm:acyclic}
If every parent has strictly smaller rank than its child, no node is its
own ancestor at any positive number of steps
\bm{BucketMath.Etymology.forest\_acyclic}.
\end{theorem}

\begin{proof}
Each step lowers the rank, so after $n+1$ steps the rank is strictly
below the start.
\end{proof}

Represent a node by its root path, the list of child indices from its
family root. Depth is the length of the path; the parent drops the last
index, so depth is a valid rank \bm{BucketMath.Etymology.address\_acyclic}.

\begin{theorem}[LCA depth]
\label{thm:lca}
The longest common prefix of two root paths is a common ancestor
\bm{BucketMath.Etymology.lcp\_common\_prefix}, and every common ancestor
has depth at most its length \bm{BucketMath.Etymology.lcp\_deepest}. Its
length is therefore the depth of the lowest common ancestor.
\end{theorem}

\begin{theorem}[Ultrametric]
\label{thm:ultra}
LCA depth satisfies $\min(s(a,b), s(b,c)) \le s(a,c)$
\bm{BucketMath.Etymology.lcaDepth\_ultra}. For any height $H$ the distance
$H - s(a,b)$ is symmetric and satisfies $d(a,c) \le \max(d(a,b), d(b,c))$
\bm{BucketMath.Etymology.rootDist\_ultrametric}.
\end{theorem}

\begin{definition}[Root distance]
\label{def:rootdist}
For lexemes $u,v$ with depths $d_u,d_v$ and LCA depth $s$, root distance
is $d_u + d_v - 2s$, the number of edges on the forest path between them.
Pairs with no common ancestor have none.
\end{definition}

\begin{theorem}[Root distance is a metric]
\label{thm:metric}
Root distance is symmetric \bm{BucketMath.Etymology.pathDist\_symm}, zero
on equal nodes \bm{BucketMath.Etymology.pathDist\_self}, zero only on
equal nodes \bm{BucketMath.Etymology.pathDist\_eq\_zero}, and satisfies
the triangle inequality \bm{BucketMath.Etymology.pathDist\_triangle}. It
is not an ultrametric: a node, its parent and a sibling give distances 1,
1 and 2 \bm{BucketMath.Etymology.pathDist\_not\_ultra}.
\end{theorem}

\begin{proof}
The triangle inequality reduces to $s(a,b)+s(b,c) \le d_b + s(a,c)$. The
smaller of the two left terms is at most $s(a,c)$ by \Cref{thm:ultra},
and the larger is at most $d_b$.
\end{proof}

T1 tests this metric. The ultrametric of \Cref{thm:ultra} orders pairs by
LCA depth alone and ignores how far each word sits below the ancestor.

A borrowed word has two lineages: its inherited parent inside its own
family and its donor in another. Give each node the set of its root paths
and set similarity to the largest LCA depth over pairs of paths.

\begin{theorem}[When borrowing breaks the ultrametric]
\label{thm:borrow}
If the middle word $b$ has one root path, the ultrametric inequality holds
for any path sets of $a$ and $c$
\bm{BucketMath.Etymology.single\_address\_ultra}. With two paths it fails:
$a$ at $[0,0,0]$, $b$ at $\{[0,0,1],[1,0,1]\}$ and $c$ at $[1,0,0]$ give
$s(a,b)=s(b,c)=2$ and $s(a,c)=0$
\bm{BucketMath.Etymology.borrowing\_breaks\_ultra}, so no proof of the
inequality for general path sets exists
\bm{BucketMath.Etymology.borrowing\_needs\_two\_addresses}.
\end{theorem}

\begin{proof}
Take $x$ in $a$'s set attaining $s(a,b)$ and $z$ in $c$'s set attaining
$s(b,c)$. Then $\min(\mathrm{lcp}(x,b),\mathrm{lcp}(b,z)) \le
\mathrm{lcp}(x,z) \le s(a,c)$. The counterexample is checked by
\texttt{decide}.
\end{proof}

The sequencing result that every acquisition order costs the same
\bm{BucketMath.Learning.Sequencing.totalCost\_perm} therefore applies to
vocabulary ordered by the inheritance forest, and stops applying as soon
as a word can be reached through a borrowing edge.

\section{Preregistered Tests}
\label{sec:tests}

Hypotheses, sampling, corrections and thresholds were committed in
\texttt{PREREGISTRATION.md} before any test statistic was computed.
Amendment 1, dated 2026-11-25 and extended 2026-11-26, changed the graph
and the T2 classes after the first run and before the second; hypotheses,
samples, alpha and corrections are unchanged. Every number in T1 to T3
comes from the second build \emp{\testrun}.

\paragraph{T1.} Out of 200,000 sampled translation pairs, 645 share a
forest ancestor, up from 333 in the first build. Among the 464 with IPA on
both sides, root distance correlates with phonetic similarity at
$\rho=-0.27$, cluster bootstrap CI $[-0.36,-0.18]$, Holm
$p=4.7\times10^{-9}$, so we reject H0a. The first build gave
$\rho=-0.41$ on 236 pairs. Related pairs average phonetic similarity 0.46
against 0.16 for the 43,609 unrelated pairs, rank-biserial $r=0.67$.
Semantic similarity has 80 related pairs with vectors, $\rho=-0.07$, CI
$[-0.30,0.15]$, Holm $p=0.48$: we fail to reject H0b. Related pairs have
higher LaBSE cosine than unrelated ones, 0.74 against 0.66
(exploratory).

\begin{figure}[htbp]
  \centering
  \includegraphics[width=0.7\linewidth]{figures/fig_distance.png}
  \caption{Mean similarity of translation pairs by root distance. Dotted
    lines: pairs with no common ancestor.}
  \label{fig:distance}
\end{figure}

\paragraph{T2.} Inherited means inheritance or derivation from the
ancestral line; borrowed means borrowing, learned borrowing or calque;
\texttt{der\_ext} is reported apart. Every per-language number depends on
this classification. Among words that translate the 1,000
most-translated English senses, 206 languages have at least 100
classified words. Median borrowed share is 0.18, and 0.23 when
\texttt{der\_ext} counts as borrowed. Frequent vocabulary has a lower
borrowed share than the whole lexicon in 193 languages; 158 reject H0 at
BH $q<0.05$, and the Fisher-pooled $p$ is below the $10^{-300}$ floor.
German goes from 0.52 over the lexicon to 0.15 in frequent words, Spanish
from 0.64 to 0.34, Japanese from 0.94 to 0.53. Spanish frequent words
split 0.30 inherited, 0.05 derived on the line, 0.18 borrowed, 0.03
\texttt{der\_ext}, 0.11 affixal and 0.33 unclassified; class shares for
every language are in \texttt{data/tests.json}. By family, median borrowed
share is 0.11 over 69 Indo-European languages and 0.29 over 16 Turkic
ones.

\begin{figure}[htbp]
  \centering
  \includegraphics[width=0.62\linewidth]{figures/fig_borrowed.png}
  \caption{Borrowed share of frequent vocabulary with Wilson 95\% CI,
    against the whole lexicon, for the 40 languages with most classified
    frequent words.}
  \label{fig:borrowed}
\end{figure}

\paragraph{T3.} Over the 23 languages of the cognate-coverage study,
genealogical closeness to English explains $R^2=0.60$ of cognate
coverage. The residuals correlate with borrowed share at $\rho=0.26$, CI
$[-0.18,0.62]$, Holm $p=0.48$. We fail to reject H0; the CI reaches past
0.4, so by the preregistered rule the result is inconclusive. Residuals
against the share borrowed from Latin, French, English or Greek give
$\rho=0.25$, $p=0.25$ (exploratory).

\begin{figure}[htbp]
  \centering
  \includegraphics[width=0.55\linewidth]{figures/fig_residuals.png}
  \caption{Cognate-coverage residual after genealogy, against borrowed
    share of frequent vocabulary.}
  \label{fig:residuals}
\end{figure}

\section{Limitations}
\label{sec:limitations}

Nodes are keyed by form, so homographs share one node and the first
descent claim wins. Only 0.32\% of sampled translation pairs reach a
common ancestor: most pairs join languages from different families, and
most translation-table words have no page whose chain reaches a shared
form. T1 rests on 464 pairs from well-documented families. The cognate
sample is ours and small, five concepts per family. Semantic similarity
uses the 27-language photon index. The classification uses the first
etymology relation on a page, and a third of Spanish frequent words have
none. English is absent from T2 because the translation tables run from
English. Genealogical closeness in T3 is near zero outside Germanic.

\paragraph{Deviations.} 2026-11-25: the first T2 run mapped a Fisher
$p$ of 0 to 1 in the Holm family through a falsy check; fixed, and
per-language $p$ clipped to $10^{-300}$ before pooling. 2026-11-25 and
2026-11-26: Amendment 1, after a review found 85\% of Spanish translation
nodes without a parent; the graph and T2 classes changed as listed in
\Cref{sec:graph} and the amendment.

\paragraph{Data and license.} Wiktionary content via Kaikki.org,
CC BY-SA 4.0 \autocite{kaikki2026raw}. The dump and graph live under
\texttt{\_intake/} and stay out of git; \texttt{make graph} and
\texttt{make data} rebuild them.

\printbibliography

\appendix

\section{Glossary}
\label{app:glossary}

\printglossaryentry{forest}
\printglossaryentry{rootdist}
\printglossaryentry{bshare}

\section{Lean Listings}
\label{app:lean}

From \texttt{lean/BucketMath/Etymology/Forest.lean}, checked by
\texttt{make lean} with Lean 4 \autocite{demoura2021lean4}:

\begin{lstlisting}[style=bucketlean]
theorem single_address_ultra (A C : List (List Nat)) (b : List Nat) :
    min (setSim A [b]) (setSim [b] C) <= setSim A C

theorem borrowing_breaks_ultra :
    setSim inherited otherFamily <
      min (setSim inherited borrowed) (setSim borrowed otherFamily) := by
  decide
\end{lstlisting}

\end{document}
